\documentclass[10pt,a4paper]{amsart}
\usepackage[margin=19mm]{geometry}
\usepackage{amsmath,amssymb,mathtools}
\usepackage[scr=rsfs,cal=euler]{mathalfa}
\usepackage{xfrac}
\usepackage[unicode,hidelinks,bookmarks=false]{hyperref}
\newcommand{\ie}{i.e.\ }
\newcommand{\cf}{cf.\ }
\newcommand{\resp}{respectively\ }
\newcommand{\Subsection}{\subsection}
\providecommand{\nobreakdash}{\nobreak\hbox{-}}
\setlength{\emergencystretch}{2em}
\multlinegap0pt
\usepackage{bm}
\usepackage{bbm}
\usepackage{dsfont}
\usepackage{xspace}
\usepackage{cancel}
\usepackage{mathtools}
\usepackage{amsthm}
\makeatletter
\@ifundefined{definition}{\newtheorem{definition}{Definition}}{}
\@ifundefined{theorem}{\newtheorem{theorem}{Theorem}}{}
\makeatother
\title[Structure preservation using discrete gradients in the Vlaso]{Structure preservation using discrete gradients in the Vlasov-Poisson-Landau system}
\author{Daniel S. Finn, Joseph V. Pusztay, Matthew G. Knepley, Mark F. Adams}
\date{Source: arXiv:2602.13068v1 (CC BY 4.0)}
\begin{document}
\maketitle

\noindent\emph{Source and licence.} Structure preservation using discrete gradients in the Vlasov-Poisson-Landau system. Version arXiv:2602.13068v1 (CC BY 4.0), distributed under the Creative Commons Attribution 4.0 International licence, as recorded in the arXiv licence metadata for this exact version and shown on the abstract page https://arxiv.org/abs/2602.13068v1. Full rights notice: licenses/arxiv-2602.13068-CC-BY-4.0.txt.\par\smallskip

\noindent\textit{Editorial scope.} Counted unit: the Theorem whose source label is \texttt{None} in the article (source0.tex lines 610--616, source label \texttt{None}), delivered together with the article's own complete original proof of it (source0.tex lines 618--638, one \texttt{proof} environment of 21 lines). No mathematics is written, repaired or re-derived: statement and proof are the article's own text line for line. Required context retained uncounted: FinDimBrack (lines 407--410); DGDef (lines 490--497); GammaBar (lines 593--596); DGColDTE (lines 598--601). Every definition, display and label of those lines is retained unchanged. Omitted from this package and not redistributed: 1--406, 411--489, 498--592, 597--597, 602--609, 617--617, 639--976. Nothing inside a retained excerpt is omitted or altered: every retained body is delivered exactly as the article's lines are written, so no omission receipt of any kind accompanies this package. Citations: the retained excerpts carry no citation command at all, so no bibliography entry of the article is transcribed and no reference is redistributed. Reference adaptation: none was needed and none was made; every reference inside the delivered unit resolves inside it, so no unresolved reference or citation is produced. Printed slips kept literal: no printed word, symbol, hypothesis or number of the counted statement or of its proof was altered. Layout and class: the article's own class is \texttt{elsarticle} with packages including amssymb, amsmath, graphicx, latexsym, amsfonts, amsthm, hyperref, xargs, todonotes, xcolor; the wrapper is the lane's pinned \texttt{amsart} editorial wrapper at 10pt on A4, which restates the theorem-like environments the retained text needs and adds the article's own macro definitions that the retained text needs (none), with extra wrapper packages (bm, bbm, dsfont, xspace, cancel, mathtools). The delivered numbering is the unit's own. Assets: no retained span contains a figure environment, a graphics-inclusion command, a PDF-page inclusion, a table or a TikZ picture; no image file is copied into this package, the package declares no asset dependency and no image file is redistributed with it. Licence: version arXiv:2602.13068v1 of this article is distributed under the Creative Commons Attribution 4.0 International licence, as recorded in the arXiv licence metadata for this exact version and shown on the abstract page https://arxiv.org/abs/2602.13068v1. A full rights notice is delivered at licenses/arxiv-2602.13068-CC-BY-4.0.txt. Characters: every retained excerpt is pure ASCII, so the delivered engine, XeLaTeX with its default Latin Modern text font, reports no missing character.
\begin{equation}
(F, G)_{h}=\frac{1}{2} \sum_{p, \overline{p}} \boldsymbol{\Gamma}(F, p, \overline{p}) \cdot \mathbb{W}(p, \overline{p}) \cdot \boldsymbol{\Gamma}(G, p, \overline{p}),
\label{FinDimBrack}
\end{equation} 

\begin{definition}
  A \textit{discrete gradient} for a smooth function $A:V \rightarrow \mathbb{R}$ is a mapping $\overline{\nabla}: V \times V \rightarrow \mathbb{R}^m$ with the following properties.
\begin{enumerate}
  \item Directionality: $\overline{\nabla} A(\boldsymbol{z}_{n+1},\boldsymbol{z}_{n}) \cdot (\boldsymbol{z}_{n}-\boldsymbol{z}_{n+1}) = A(\boldsymbol{z}_{n}) - A(\boldsymbol{z}_{n+1})$ for any $\boldsymbol{z}_{n+1},\boldsymbol{z}_{n} \in V$.
  \item Consistency: $\overline{\nabla} A(\boldsymbol{z}_{n+1},\boldsymbol{z}_{n})=\nabla A(\frac{\boldsymbol{z}_{n+1}+\boldsymbol{z}_{n}}{2}) + \mathcal{O}(||\boldsymbol{z}_{n}-\boldsymbol{z}_{n+1}||)$ for all $\boldsymbol{z}_{n+1},\boldsymbol{z}_{n} \in V$ with $||\boldsymbol{z}_{n}-\boldsymbol{z}_{n+1}||$ sufficiently small.
  \end{enumerate}
  \label{DGDef}
\end{definition}

\begin{equation}
\overline{\boldsymbol{\Gamma}_{n}^{n+1}(F, p, \overline{p})}=\frac{1}{m w_{p}} \overline{\nabla} F (\boldsymbol{v}_{p}^{n},\boldsymbol{v}_{p}^{n+1})-\frac{1}{m w_{\overline{p}}} \overline{\nabla} F (\boldsymbol{v}_{\overline{p}}^{n},\boldsymbol{v}_{\overline{p}}^{n+1}).
\label{GammaBar}
\end{equation}

\begin{equation}
\frac{F^{n+1} - F^{n}}{\Delta t} = \frac{1}{2} \sum_{p,\overline{p}} \overline{\boldsymbol{\Gamma}^{n+1}_{n} (F,p,\overline{p})} \cdot \overline{W^{n+1}_{n} (p,\overline{p})} \cdot \overline{\boldsymbol{\Gamma}^{n+1}_{n} (S_{\epsilon},p,\overline{p})},
\label{DGColDTE}
\end{equation}

\begin{theorem}
The collisional discrete-time evolution equation, given by~\eqref{DGColDTE}, is a consistent approximation to the finite-dimensional bracket~\eqref{FinDimBrack}. \\
\begin{equation}
\left. \frac{d F}{d t}\right|_{\text {coll }} \approx \frac{F^{n+1} - F^{n}}{\Delta t},
\end{equation}
in the limit where $||\boldsymbol{v}_{p}^{n+1} - \boldsymbol{v}_{p}^{n}|| \rightarrow 0$ and $||\boldsymbol{v}_{\overline{p}}^{n+1} - \boldsymbol{v}_{\overline{p}}^{n}|| \rightarrow 0$.
\end{theorem}

\begin{proof}
First we consider what happens to~\eqref{GammaBar} as $||\boldsymbol{v}_{p}^{n+1} - \boldsymbol{v}_{p}^{n}||$ and $||\boldsymbol{v}_{\overline{p}}^{n+1} - \boldsymbol{v}_{\overline{p}}^{n}||$ approach zero, 
\begin{align*}
\overline{\boldsymbol{\Gamma}_{n}^{n+1}(F, p, \overline{p})}=&\frac{1}{m w_{p}} \overline{\nabla} F (\boldsymbol{v}_{p}^{n},\boldsymbol{v}_{p}^{n+1})-\frac{1}{m w_{\overline{p}}} \overline{\nabla} F (\boldsymbol{v}_{\overline{p}}^{n},\boldsymbol{v}_{\overline{p}}^{n+1})\\
= &\frac{1}{m w_{p}} \left( \nabla F \left( \boldsymbol{v}_{p}^{n+1/2} \right) + \mathcal{O} (||\boldsymbol{v}_{p}^{n+1} - \boldsymbol{v}_{p}^{n}||) \right) \\
&- \frac{1}{m w_{\overline{p}}} \left( \nabla F \left( \boldsymbol{v}_{\overline{p}}^{n+1/2} \right) + \mathcal{O} (||\boldsymbol{v}_{\overline{p}}^{n+1} - \boldsymbol{v}_{\overline{p}}^{n}||) \right)\\
\approx &\frac{1}{m w_{p}} \nabla F \left( \boldsymbol{v}_{p}^{n+1/2} \right) - \frac{1}{m w_{\overline{p}}} \nabla F \left( \boldsymbol{v}_{\overline{p}}^{n+1/2} \right)\\
\approx & \boldsymbol{\Gamma} (F, p, \overline{p})
\end{align*}
where $\boldsymbol{v}_p^{n+1/2}=\frac{1}{2}\left(\boldsymbol{v}_p^{n+1}+\boldsymbol{v}_p^{n}\right)$ and likewise for $\boldsymbol{v}_{\bar{p}}^{n+1/2}$. Furthermore, we recall the consistency property in Definition~\eqref{DGDef} to obtain the midpoint approximation, $\nabla F(\boldsymbol{v}_p^{n+1/2})$, of the discrete gradient, $\overline{\nabla} F (\boldsymbol{v}_{p}^{n},\boldsymbol{v}_{p}^{n+1})$. Plugging this into relation into~\eqref{DGColDTE} gives,
\begin{align*}
\frac{F^{n+1} - F^{n}}{\Delta t} &= \frac{1}{2} \sum_{p,\overline{p}} \overline{\boldsymbol{\Gamma}^{n+1}_{n} (F,p,\overline{p})} \cdot \overline{W^{n+1}_{n} (p,\overline{p})} \cdot \overline{\boldsymbol{\Gamma}^{n+1}_{n} (S_{\epsilon},p,\overline{p})} \\
&\approx \frac{1}{2} \sum_{p,\overline{p}} \boldsymbol{\Gamma}^{n+1}_{n} (F,p,\overline{p}) \cdot \left[ \nu w_{p} w_{\overline{p}} \boldsymbol{1}(p, \overline{p}) \mathbb{Q}\left(\boldsymbol{\Gamma}_{n}^{n+1}(K, p, \overline{p})\right) \right] \cdot \boldsymbol{\Gamma}^{n+1}_{n} (S_{\epsilon},p,\overline{p})
\end{align*}
Realizing $\boldsymbol{\Gamma}(K,p,\overline{p}) = \boldsymbol{v}_p - v_{\overline{p}}$, we can therefore show that,
\begin{align*}
\frac{F^{n+1} - F^{n}}{\Delta t} &\approx \frac{1}{2} \sum_{p,\overline{p}} \boldsymbol{\Gamma}^{n+1}_{n} (F,p,\overline{p}) \cdot \left[ \nu w_{p} w_{\overline{p}} \boldsymbol{1}(p, \overline{p}) \mathbb{Q}\left( \boldsymbol{v}_p - v_{\overline{p}}\right) \right] \cdot \boldsymbol{\Gamma}^{n+1}_{n} (S_{\epsilon},p,\overline{p}) \\
&\approx \frac{1}{2} \sum_{p,\overline{p}} \boldsymbol{\Gamma}^{n+1}_{n} (F,p,\overline{p}) \cdot \mathbb{W}(p, \overline{p}) \cdot \boldsymbol{\Gamma}^{n+1}_{n} (S_{\epsilon},p,\overline{p})\\
&\approx  \left. \frac{d F}{d t}\right|_{\text {coll }}
\end{align*}
\end{proof}

\end{document}
